\section{Continual Learning、World Model 与 Expert Agent}

当主持人问 Agent 最大瓶颈是什么，苏煜把 memory、self learning、continual learning、world model、specialization、expert agent 统一成同一件事的不同侧面。Agent 现在缺的，是从交互中持续学习，并把经验变成稳定能力。

\begin{figure}[H]
\centering
\includegraphics[width=0.94\textwidth]{continual-learning-map.png}
\caption{Continual Learning、World Model 与 Expert Agent 的关系。}
\end{figure}

\begin{knowledgebox}{读图：这些术语为什么是一条链}
左侧的经验来自环境交互；continual/self learning 把经验更新成能力；world model 表示对环境规律的学习；specialization 表示形成领域专家；最后结果是更可靠、更快、更低成本的 Agent。它支持的结论是：memory、world model、专家化不是彼此独立的 roadmap，而是同一条能力积累链。
\end{knowledgebox}

\subsection{术语消化：Agent 瓶颈词表}

\noindent\begin{tabularx}{\textwidth}{p{0.20\textwidth}p{0.32\textwidth}X}
\toprule
术语 & 解决的问题 & 本期中的位置 \\
\midrule
Memory & 多步任务中保留相关历史与偏好 & 没有记忆，Agent 每次都像新手，容易重复犯错。 \\
Continual Learning & 从新交互中持续更新能力 & 让 Agent 不只是执行脚本，而是积累经验。 \\
World Model & 学到环境状态、因果与可行动性 & 支撑规划、预测后果和安全操作。 \\
Specialization & 成为某领域 expert agent & 从通用助手变成可靠的垂直工作者。 \\
Reliability & 任务稳定完成率 & 当前 Agent 产品化最大痛点之一。 \\
Cost Effectiveness & 单任务成本是否可接受 & 决定 Agent 能否从 demo 进入生产。 \\
\bottomrule
\end{tabularx}

\begin{importantbox}{最大瓶颈的统一表述}
Agent 最大瓶颈不是单个 memory 模块、单个 prompt 或单个工具，而是无法像人一样把长期经验转化成稳定专业能力。Continual learning、world model、specialization 最终都服务于可靠性、速度和成本。
\end{importantbox}

\subsection{Forward-deployed engineers 的信号}

访谈提到，一些公司采用 forward-deployed engineers，派工程师到客户现场帮他们 build agent。这说明当前 Agent 还没有低门槛普及：需要理解客户流程、工具环境、失败模式和安全约束。产品仍处在高服务含量阶段。

\begin{warningbox}{部署难不只是模型弱}
Agent 进企业难，不只是模型能力不足，也包括流程理解、权限管理、数据接入、错误责任、安全边界和成本核算。把所有失败都归因于“模型不够聪明”会误判产品化难度。
\end{warningbox}

\subsection{本章小结}

Agent 的下一阶段主线是持续学习与世界建模。真正的进步会表现为可靠、快速、低成本和专家化，而不是只在 demo 中完成更炫的动作序列。

